% Part: first-order-logic
% Chapter: completeness
% Section: lindenbaums-lemma

\documentclass[../../../include/open-logic-section]{subfiles}

\begin{document}

\iftag{FOL}
      {\olfileid{fol}{com}{lin}}
      {\olfileid{pl}{com}{lin}}

\olsection{لم لیندنباوم}

\begin{explain}
اکنون لمی را اثبات می‌کنیم که نشان می‌دهد هر مجموعهٔ سازگار از
!!{sentence}s در مجموعه‌ای از جمله‌ها گنجانده می‌شود که نه‌فقط سازگار،
بلکه !!{complete} نیز هست. اثبات با افزودن هر بار یک !!{sentence} پیش
می‌رود و در هر گام تضمین می‌کند که مجموعه سازگار بماند. این کار را چنان
انجام می‌دهیم که برای هر~$!A$، در مرحله‌ای یا $!A$ افزوده شود یا
$\lnot !A$. آنگاه اجتماع همهٔ مراحل این ساخت یا $!A$ را دربر دارد یا
نقیض آن~$\lnot !A$ را، و ازاین‌رو کامل است. این اجتماع همچنین سازگار
است، زیرا در هر مرحله مراقبیم ناسازگاری‌ای پدید نیاوریم.
\end{explain}

\begin{lem}[لم لیندنباوم]
\ollabel{lem:lindenbaum} هر مجموعهٔ سازگار~$\Gamma$ در زبان~$\Lang{L}$
را می‌توان به !!a{complete} مجموعهٔ سازگار~$\Gamma^*$ گسترش داد.
\end{lem}

\begin{proof}
فرض کنید $\Gamma$ سازگار باشد. بگذارید $!A_0$، $!A_1$،
\dots{} شمارشی از همهٔ !!{sentence}s زبان~$\Lang L$ باشد.
$\Gamma_0 = \Gamma$ را تعریف کنید، و
\[
\Gamma_{n+1} =
\begin{cases}
\Gamma_n \cup \{ !A_n \} & \textrm{اگر $\Gamma_n \cup \{!A_n\}$
  سازگار باشد؛} \\
\Gamma_n \cup \{ \lnot !A_n \} & \textrm{در غیر این صورت.}
\end{cases}
\]
بگذارید $\Gamma^* = \bigcup_{n \geq 0} \Gamma_n$.

هر $\Gamma_n$ سازگار است: $\Gamma_0$ بنا بر تعریف سازگار است. اگر
$\Gamma_{n+1} = \Gamma_n \cup \{!A_n\}$، علتش آن است که مجموعهٔ اخیر
سازگار است. اگر چنین نباشد، $\Gamma_{n+1} = \Gamma_n \cup \{\lnot
!A_n\}$. باید بررسی کنیم که $\Gamma_n \cup \{\lnot !A_n\}$ سازگار
است. فرض کنید سازگار نباشد. آنگاه \emph{هر دو} مجموعهٔ
$\Gamma_n \cup \{!A_n\}$ و $\Gamma_n \cup \{\lnot !A_n\}$ ناسازگار
هستند. این بدان معناست که بنا بر
\iftag{FOL}{%
  \tagrefs{prfAX/{fol:axd:prv:prop:provability-exhaustive},
    prfSC/{fol:seq:prv:prop:provability-exhaustive},
    prfND/{fol:ntd:prv:prop:provability-exhaustive},
    prfTab/{fol:tab:prv:prop:provability-exhaustive}}}{%
  \tagrefs{prfAX/{pl:axd:prv:prop:provability-exhaustive},
    prfSC/{pl:seq:prv:prop:provability-exhaustive},
    prfND/{pl:ntd:prv:prop:provability-exhaustive},
    prfTab/{pl:tab:prv:prop:provability-exhaustive}}}،
$\Gamma_n$ ناسازگار می‌بود، برخلاف فرض استقرا.

برای هر~$n$ و هر $i < n$، $\Gamma_i \subseteq \Gamma_n$. این امر با
استقرایی ساده بر~$n$ نتیجه می‌شود. برای $n=0$ هیچ $i < 0$ای وجود ندارد،
پس ادعا خودبه‌خود برقرار است. برای گام استقرا، فرض کنید ادعا برای~$n$
صادق باشد. نشان می‌دهیم اگر $i < n+1$، آنگاه
$\Gamma_i \subseteq \Gamma_{n+1}$. بنا بر ساخت، یا
$\Gamma_{n+1} = \Gamma_n \cup \{!A_n\}$ یا
$= \Gamma_n \cup \{\lnot !A_n\}$. پس
$\Gamma_n \subseteq \Gamma_{n+1}$. اگر $i < n+1$، آنگاه بنا بر فرض
استقرا $\Gamma_i \subseteq \Gamma_n$ (اگر $i < n$)، یا بنا بر امر بدیهی
$\Gamma_n \subseteq \Gamma_n$ (اگر $i = n$). از تعدی~$\subseteq$
نتیجه می‌گیریم که $\Gamma_i \subseteq \Gamma_{n+1}$.

از این امر نتیجه می‌شود که $\Gamma^*$ سازگار است. دلیل چنین است:
بگذارید $\Gamma' \subseteq \Gamma^*$ متناهی باشد. اگر این زیرمجموعه
تهی باشد، سازگار است. در غیر این صورت،
هر $!B \in \Gamma'$ برای عددی~$i$ در~$\Gamma_i$ نیز هست. بگذارید $n$
بزرگ‌ترین این اعداد باشد. چون اگر $i \le n$، آنگاه
$\Gamma_i \subseteq \Gamma_n$، هر $!B \in \Gamma'$ همچنین
$\in \Gamma_n$ است؛ یعنی $\Gamma' \subseteq \Gamma_n$، و $\Gamma_n$
سازگار است. پس هر
زیرمجموعهٔ متناهی $\Gamma' \subseteq \Gamma^*$ سازگار است. بنا بر
\iftag{FOL}{%
  \tagrefs{prfAX/{fol:axd:ptn:prop:proves-compact},
    prfSC/{fol:seq:ptn:prop:proves-compact},
    prfND/{fol:ntd:ptn:prop:proves-compact},
    prfTab/{fol:tab:ptn:prop:proves-compact}}}{%
  \tagrefs{prfAX/{pl:axd:ptn:prop:proves-compact},
    prfSC/{pl:seq:ptn:prop:proves-compact},
    prfND/{pl:ntd:ptn:prop:proves-compact},
    prfTab/{pl:tab:ptn:prop:proves-compact}}}، $\Gamma^*$ سازگار است.

هر !!{sentence} در~$\Frm[L]$ در فهرستی ظاهر می‌شود که برای
تعریف~$\Gamma^*$ به کار رفت. اگر $!A_n \notin \Gamma^*$، علتش آن است
که $\Gamma_n \cup \{!A_n\}$ ناسازگار بوده است. اما آنگاه
$\lnot !A_n \in \Gamma^*$؛ پس $\Gamma^*$ !!{complete} است.
\end{proof}

\end{document}
